% Auto-generated — REAL cluster timings (mode=cluster).
% Highlight: N=20 / C5  runs=3

\subsection{Orthogonal Testing Efficiency}

\begin{figure*}[t]
\centering
% \includegraphics[width=0.85\linewidth]{figures/evaluation/orthogonal-reduction.png}
\caption{Orthogonal validation under different scales (real cluster).
(a) Test-case growth vs configuration items.
(b) Complete validation wall-clock time vs configuration items.}
\label{fig:orthogonal-reduction}
\end{figure*}

\textbf{Goal.}
Evaluate orthogonal reduction of dynamic validation cost on a live Istio cluster, and whether fault detection coverage is preserved after merging.

\textbf{Setup.}
Service scales $N\in\{15,20\}$, configuration levels $\{C4,C5\}$.
Each validation case performs live \texttt{kubectl apply}, ingress traffic, assertion, and delete.
Orthogonal bundles share one apply/clear and issue per-effect traffic.
Each cell is repeated 3 times; times are wall-clock mean$\pm$std.

\textbf{Results.}
Under a 20-service deployment at C5 (590 configuration items), exhaustive validation uses 400 cases vs 85 orthogonal cases (reduction 78.75\%).
Complete validation wall-clock time drops from $316.4\pm0.2$s to $79.6\pm0.3$s (time reduction 74.84\%).

\textbf{Coverage ablation.}
Table~\ref{tab:ortho-ablation} reports live fault/effect detection before vs after merging.

\begin{table*}[t]
\centering
\caption{Real-cluster detection and validation time (C5).}
\label{tab:ortho-ablation}
\scriptsize
\begin{tabular}{cccccccccc}
\toprule
$N$ & $|F|$ & DR$_b$ & DR$_a$ & $\Delta$ & Cases$_e$ & Cases$_o$ & Red. & $T_e$ (s) & $T_o$ (s) \\
\midrule
15 & 300 & 100.00\% & 100.00\% & 0.00 & 300 & 65 & 78.33\% & 234.4$\pm$0.2 & 58.4$\pm$0.1 \\
20 & 400 & 100.00\% & 100.00\% & 0.00 & 400 & 85 & 78.75\% & 316.4$\pm$0.2 & 79.6$\pm$0.3 \\
\bottomrule
\end{tabular}
\end{table*}
